\documentclass[10pt,twocolumn]{article}

\usepackage[margin=0.75in]{geometry}
\usepackage{amsmath}
\usepackage{amssymb}
\usepackage{booktabs}
\usepackage{array}
\usepackage{enumitem}
\usepackage{graphicx}
\usepackage[hidelinks]{hyperref}
\usepackage{xcolor}

\setlist[itemize]{leftmargin=*,itemsep=1pt,topsep=2pt}
\setlist[enumerate]{leftmargin=*,itemsep=1pt,topsep=2pt}
\renewcommand{\arraystretch}{1.12}
\sloppy

\title{\vspace{-0.5em}AI Agents for Love Letter:\\
Reasoning, Belief Search, and Reinforcement Learning under Hidden Information}

\author{
Yanqi Zhang \and Yueqian Ma \and Yunyang Liu\\
CS181 Final Project
}

\date{}

\begin{document}
\maketitle

\begin{abstract}
This project studies \emph{Love Letter} as a compact stochastic,
partially observable, adversarial game. We implement a reusable Python game
environment and compare several agent families from the CS181 course:
random play, handcrafted heuristics, symbolic logic, approximate Bayesian
belief search, expectimax, Monte Carlo tree search, tabular Q-learning, and
feature-based approximate Q-learning.
The central implementation question is not simply whether an agent can win
against a weak baseline, but whether it does so under the same hidden
information available to a real player. Therefore, the project separates
fair hidden-information agents from perfect-information oracle baselines.
Pilot experiments show that card-counting heuristics, logic constraints, and
Q-learning with heuristic priors and reward-shaped approximate Q-learning
consistently outperform random play, while belief search is harder to tune and
has high variance at small sample sizes.
The result is a complete experimental platform for studying uncertainty,
search, and learning in a small but nontrivial imperfect-information game.
\end{abstract}

\section{Introduction}

Love Letter is a small card game with only 16 cards, but it contains the
features that make many AI problems difficult: private hands, a hidden
burned card, stochastic draws, adversarial targeting, and actions that
reveal partial information. This makes it a useful course project because
the full rules are implementable within a semester, while the decision
problem still requires reasoning under uncertainty.

Our proposal framed the project around four agent families: heuristic
agents, Bayesian agents, expectimax search, and reinforcement learning with
MCTS. The current implementation follows that plan and refactors the branch
work into a single package. In addition, we explicitly distinguish two kinds
of search:
\begin{itemize}
    \item \textbf{Fair hidden-information search}, where the agent samples
    hidden states from an observation-derived belief and does not directly
    inspect the opponent's true hand.
    \item \textbf{Oracle search baselines}, where the agent searches from a
    cloned true environment. These baselines measure an upper bound for
    planning strength, but they are not directly comparable to real
    hidden-information play.
\end{itemize}

This distinction is important for soundness. If an agent sees the full
environment state, high win rates are less meaningful because they may come
from illegal information. The report therefore presents oracle agents as
diagnostic baselines rather than as the main claim.

\section{Game Model}

The base deck is shown in Table~\ref{tab:deck}. Each turn consists of drawing
one card, then playing one of the two cards in hand. A round ends when all but
one player are eliminated or when the deck is exhausted and the surviving
players compare their remaining cards.

\begin{table}[t]
\centering
\caption{Base Love Letter deck used by the environment.}
\label{tab:deck}
\begin{tabular}{lrrp{0.36\columnwidth}}
\toprule
Card & Value & Count & Effect summary\\
\midrule
Guard & 1 & 5 & Guess a non-Guard card.\\
Priest & 2 & 2 & Privately view a hand.\\
Baron & 3 & 2 & Compare hands; lower card loses.\\
Handmaid & 4 & 2 & Gain temporary protection.\\
Prince & 5 & 2 & Force discard and redraw.\\
King & 6 & 1 & Swap hands.\\
Countess & 7 & 1 & Forced with King or Prince.\\
Princess & 8 & 1 & Lose if discarded.\\
\bottomrule
\end{tabular}
\end{table}

We model a complete state as
\begin{equation}
s = (s_{\mathrm{pub}}, s_{\mathrm{priv}}),
\end{equation}
where $s_{\mathrm{pub}}$ includes discard piles, eliminated players,
protection flags, turn index, current player, and deck size, while
$s_{\mathrm{priv}}$ includes private hands, the hidden set-aside card, and
the exact deck order. For player $i$, the observation function is
\begin{equation}
o_i = O_i(s) =
(\mathrm{public\ view}, \mathrm{own\ hand}, \mathrm{private\ events}_i).
\end{equation}
Agents receive $o_i$ and the legal action set $A(o_i)$.

The transition model is mixed:
\begin{itemize}
    \item Card effects such as Guard, Baron, King, and Princess are
    deterministic once the hidden state is fixed.
    \item Draws are stochastic from the agent's perspective because the deck
    order and hidden card are unknown.
\end{itemize}
For reinforcement learning, terminal reward is implemented as $+1$ for a win
and $-1$ for a loss. This differs slightly from the proposal's $0$ loss
reward, but it provides a clearer learning signal and separates losing
terminal states from neutral nonterminal transitions.

\section{Environment Implementation}

The environment is implemented in \texttt{loveletter\_ai.environment}.
\texttt{LoveLetterEnv} owns the true state, legal action generation, seeded
randomness, and card-effect resolution. The true state is represented by
\texttt{GameState}; the hidden-information-safe player view is represented by
\texttt{Observation}. This boundary is the main engineering device that keeps
agents fair.

Important implementation choices include:
\begin{itemize}
    \item \textbf{Deterministic reproducibility.} Each round uses an explicit
    seed. This makes tournament and ablation results repeatable.
    \item \textbf{Rule-complete legal actions.} Legal action generation handles
    Countess forced play, Guard guess restrictions, Handmaid protection, and
    Prince self-targeting.
    \item \textbf{Public/private event history.} Priest observations are
    private to the viewer, while public events such as discards and
    eliminations are visible to all agents.
    \item \textbf{Search cloning.} The environment supports deep cloning for
    expectimax and MCTS rollouts. The fair search agents overwrite hidden
    fields with belief samples before rollouts.
\end{itemize}

The evaluation layer in \texttt{loveletter\_ai.evaluation} provides
\texttt{run\_round}, \texttt{evaluate\_pair}, and
\texttt{evaluate\_pair\_symmetric}. Symmetric evaluation runs both seat
assignments, reducing first-player and seed bias.

\section{Agents}

\subsection{Random and Heuristic Agents}

The random agent samples uniformly from legal actions and serves as the
baseline. The naive heuristic agent scores actions with simple rules such as
protecting the Princess, preferring Guard guesses, and avoiding Baron when the
retained card is weak.

The improved heuristic and greedy agents add card-count awareness. From an
observation $o$, the agent estimates unseen card multiplicities
\begin{equation}
C_o(c) = C_{\mathrm{deck}}(c) - C_{\mathrm{own}}(c) -
C_{\mathrm{discard}}(c).
\end{equation}
This lets it avoid impossible Guard guesses and estimate opponent hand values.
The action score has the general form
\begin{equation}
\begin{aligned}
H(o,a) ={}& w_r \cdot \mathrm{retained}(o,a)
+ w_t \cdot \mathrm{target}(o,a)\\
&+ w_g \cdot \mathrm{guess}(o,a)
+ w_s \cdot \mathrm{safety}(o,a).
\end{aligned}
\end{equation}
The weights are handcrafted in code rather than learned. This makes the policy
easy to inspect and provides a strong baseline for the more advanced agents.

\subsection{Logic Agent}

The logic agent maintains a persistent knowledge base of possible opponent
cards. The set for opponent $j$ is initialized as all card types and updated
by public and private evidence:
\begin{itemize}
    \item visible own cards and discard piles are removed;
    \item Priest observations collapse a target's possible set to one card;
    \item failed Guard guesses remove the guessed card;
    \item King swaps reset affected hand facts because ownership changes.
\end{itemize}
If the agent holds a Guard and the knowledge base has a singleton target set,
it makes the guaranteed guess. Otherwise it falls back to the improved
heuristic policy.

\subsection{Belief Expectimax}

The Bayesian component builds an approximate posterior over hidden states from
the observation and logic constraints. It is approximate because it does not
enumerate all histories or model opponent policies. Instead, it samples
opponent hands, the burned card, and deck order without replacement from
remaining card counts:
\begin{equation}
b_i(h \mid o_i) \propto
\mathbf{1}[h \text{ is consistent with } o_i]\prod_{c} C_o(c).
\end{equation}

Belief expectimax evaluates a root action by sampling hidden states
$h^{(1)},\ldots,h^{(K)}$ and averaging the expectimax value:
\begin{equation}
\hat{Q}(o_i,a) = \frac{1}{K}\sum_{k=1}^{K}
V_{\mathrm{exp}}(s_{\mathrm{pub}}, h^{(k)}, a).
\end{equation}
Opponent actions are treated as chance nodes with a uniform policy. Draw nodes
are enumerated by card multiplicity when possible.

\subsection{MCTS}

The oracle MCTS agent uses UCT over cloned complete states:
\begin{equation}
\mathrm{UCT}(v') =
\frac{Q(v')}{N(v')} +
c\sqrt{\frac{\log N(v)}{N(v')}}.
\end{equation}
Rollouts use an improved heuristic for the root player and random play for
opponents.

The fair \texttt{BeliefMCTSAgent} performs root information-set search. Each
simulation samples a plausible hidden state from the belief, evaluates one
root action, and updates root-level action statistics. This is less powerful
than a full information-set MCTS tree, but it avoids directly reading hidden
opponent cards and is therefore aligned with the project setting.

\subsection{Q-Learning}

The first Q-learning agent uses tabular state-action values over compact
observation and action keys. The update rule is
\begin{equation}
Q(o,a) \leftarrow Q(o,a) +
\alpha\left[r + \gamma\max_{a'}Q(o',a') - Q(o,a)\right].
\end{equation}
Because the Love Letter observation space is sparse, the agent combines learned
Q-values with a heuristic prior:
\begin{equation}
\mathrm{score}(o,a) = Q(o,a) + \lambda H_{\mathrm{norm}}(o,a).
\end{equation}
The improved training loop uses linearly decaying exploration, per-state
learning-rate decay, mixed opponents, and small reward shaping terms. The
mixed opponent pool samples Random, Naive, Improved, and self-play opponents
so that the policy does not overfit to one behavior distribution. The shaping
terms give dense feedback for tactically clear events: eliminating an opponent,
being eliminated, discarding Princess, protecting Princess with Handmaid, and
making an informative Priest move.

We also implement an approximate Q-learning agent with a linear value model
\begin{equation}
Q(o,a) = \sum_k w_k f_k(o,a).
\end{equation}
The features include retained card value, normalized heuristic score, Guard
guess probability, impossible Guard guesses, Baron value edge, Prince
Princess-discard probability, and Princess safety indicators. Tactically
monotone features are constrained to keep sensible signs, for example
discarding Princess remains non-positive and protecting Princess remains
non-negative. This makes approximate Q-learning a stable residual correction
over the heuristic policy rather than a noisy replacement for it.

\section{Experiments}

\subsection{Protocol}

All reported pilot experiments use two-player rounds. For each pair, we use
symmetric evaluation: agent A plays as player 0 for $n$ seeds, then as player
1 for another $n$ seeds. Win rate is the fraction of all $2n$ rounds won by
the first listed agent. Confidence intervals are Wilson 95 percent intervals.

The current report tables were generated with:
\begin{center}
\texttt{python test\_agent/run\_ablation\_experiments.py 20 outputs/report\_experiments quick}
\end{center}
Thus each row contains 40 games. These results are suitable for development
analysis and presentation drafts, but the final submission should rerun the
full profile with a larger sample size.

\subsection{Fair Hidden-Information Results}

Table~\ref{tab:fair} compares fair hidden-information agents against Random.
The strongest simple agents are the logic, greedy, and improved heuristic
agents, each reaching 77.5 percent in this pilot run. The repaired RL agents
now match that tier: tabular Q with a heuristic prior reaches 77.5 percent,
and approximate Q with shaping also reaches 77.5 percent. BeliefMCTS improves
over the earlier belief-search result, while BeliefExpectimax remains more
sensitive to sampling noise.

\begin{table}[t]
\centering
\caption{Fair hidden-information agents vs Random, 20 games per seat.}
\label{tab:fair}
\begin{tabular}{lrrr}
\toprule
Agent & W-L-D & Win rate & 95\% CI\\
\midrule
Naive & 27-13-0 & 67.5\% & 52.0--79.9\\
Logic & 31-9-0 & 77.5\% & 62.5--87.7\\
Greedy & 31-9-0 & 77.5\% & 62.5--87.7\\
Improved & 31-9-0 & 77.5\% & 62.5--87.7\\
BeliefExpectimax & 26-14-0 & 65.0\% & 49.5--77.9\\
BeliefMCTS & 30-10-0 & 75.0\% & 59.8--85.8\\
Tabular Q+prior & 31-9-0 & 77.5\% & 62.5--87.7\\
Approx Q+shaping & 31-9-0 & 77.5\% & 62.5--87.7\\
\bottomrule
\end{tabular}
\end{table}

\subsection{Ablation Results}

Table~\ref{tab:ablations} summarizes the main ablations. Three patterns are
important. First, pure tabular Q-learning is weak because the state-action
space is sparse. Second, adding a heuristic prior recovers most of the lost
performance, and approximate Q-learning with shaping gives a small further
gain. Third, more computation does not automatically improve a
hidden-information search agent at small sample size: extra samples can
increase variance if the belief model or rollout policy is weak.

\begin{table*}[t]
\centering
\caption{Pilot ablations. Each row is evaluated against Random with 20 games per seat.}
\label{tab:ablations}
\begin{tabular}{llrrr}
\toprule
Family & Setting & W-L-D & Win rate & 95\% CI\\
\midrule
BeliefExpectimax samples & samples=1 & 30-10-0 & 75.0\% & 59.8--85.8\\
BeliefExpectimax samples & samples=4 & 29-11-0 & 72.5\% & 57.2--83.9\\
BeliefExpectimax samples & samples=8 & 30-10-0 & 75.0\% & 59.8--85.8\\
BeliefExpectimax depth & depth=1 & 27-13-0 & 67.5\% & 52.0--79.9\\
BeliefExpectimax depth & depth=2 & 27-13-0 & 67.5\% & 52.0--79.9\\
BeliefMCTS simulations & simulations=4 & 26-14-0 & 65.0\% & 49.5--77.9\\
BeliefMCTS simulations & simulations=12 & 22-18-0 & 55.0\% & 39.8--69.3\\
BeliefMCTS simulations & simulations=24 & 27-13-0 & 67.5\% & 52.0--79.9\\
Q-learning variants & heuristic only & 30-10-0 & 75.0\% & 59.8--85.8\\
Q-learning variants & pure tabular Q & 24-16-0 & 60.0\% & 44.6--73.7\\
Q-learning variants & tabular Q+prior & 30-10-0 & 75.0\% & 59.8--85.8\\
Q-learning variants & approx Q+shaping & 31-9-0 & 77.5\% & 62.5--87.7\\
Approx Q episodes & episodes=0 & 29-11-0 & 72.5\% & 57.2--83.9\\
Approx Q episodes & episodes=500 & 30-10-0 & 75.0\% & 59.8--85.8\\
Approx Q episodes & episodes=1000 & 30-10-0 & 75.0\% & 59.8--85.8\\
Approx Q episodes & episodes=5000 & 31-9-0 & 77.5\% & 62.5--87.7\\
Oracle Expectimax & depth=1 & 30-10-0 & 75.0\% & 59.8--85.8\\
Oracle Expectimax & depth=2 & 33-7-0 & 82.5\% & 68.0--91.2\\
Oracle MCTS & simulations=4 & 32-8-0 & 80.0\% & 65.2--89.5\\
Oracle MCTS & simulations=12 & 32-8-0 & 80.0\% & 65.2--89.5\\
\bottomrule
\end{tabular}
\end{table*}

\section{Analysis}

The main observation is that the fair agents do not usually exceed 80 percent
against Random in the current pilot setting. This is reasonable. Love Letter
has high inherent variance: an early Guard hit, Princess draw, or Prince
discard can decide a round before a policy's long-term advantage has time to
show. Because each row currently contains only 40 rounds, the confidence
intervals are wide.

The most reliable improvement comes from eliminating clearly bad actions:
avoiding impossible Guard guesses, respecting the Countess rule, protecting a
Princess with Handmaid, and using remaining-card counts for target decisions.
That explains why the improved heuristic, greedy, and logic agents are strong
relative to their complexity.

Belief search is conceptually more aligned with the hidden-information game,
but its current implementation is still approximate. It samples plausible
states, assumes simple opponent behavior, and uses relatively shallow search.
When the belief is noisy, deeper search can amplify a bad determinization
rather than fix it. This explains why the sample and depth ablations do not
yet show monotonic improvement.

The Q-learning ablation confirms the diagnosis from development. Pure tabular
Q-learning is not enough: it reaches only 60 percent in the pilot table.
Adding the heuristic prior recovers the baseline to 75 percent, while the
approximate Q-learning agent with reward shaping reaches 77.5 percent. This
does not prove a large statistically significant RL advantage yet, because the
confidence intervals are wide, but it gives a much cleaner story: each repair
addresses a concrete weakness of the previous RL version.

\section{Current Shortcomings}

Compared with the reference report, this project is strongest in environment
soundness and hidden-information framing, but weaker in final experimental
scale. The main remaining issues are:
\begin{itemize}
    \item \textbf{Sample size.} The current report tables use pilot-scale
    experiments. A final report should use the full ablation profile and a
    larger symmetric tournament.
    \item \textbf{Tournament matrix.} The proposal promised a full agent
    tournament. The code supports pairwise evaluation, but the final report
    should include an all-pairs matrix similar to the reference report's Sente
    and Gote tables.
    \item \textbf{Belief model.} The Bayesian agent is an approximate
    consistent-state sampler, not a full posterior over histories.
    \item \textbf{RL maturity.} The repaired RL agent now has meaningful
    ablations, but the final report should still include larger learning
    curves before claiming a statistically significant advantage.
    \item \textbf{Two-player focus.} The engine supports 2--4 players, but the
    reported experiments use two-player rounds for controlled comparison.
\end{itemize}

\section{External Resources and Tools}

No external Love Letter engine code was copied. The game engine, agents, and
evaluation scripts are implemented in Python within this repository. External
resources and tools used are:
\begin{itemize}
    \item \textbf{Python standard library}: dataclasses, enum, random, copy,
    collections, and pathlib are used for state structures, seeded stochastic
    transitions, cloning, and counters.
    \item \textbf{pytest}: used for unit and integration tests covering setup,
    hidden-information rendering, belief sampling, agent legality, and match
    completion.
    \item \textbf{Tectonic/LaTeX}: used to compile this report from the LaTeX
    source.
    \item \textbf{Codex AI coding assistant}: used as a development assistant
    for branch integration, refactoring, testing, report drafting, and
    experiment organization. The team should review all code and explanations
    before final submission.
    \item \textbf{Love Letter base rules}: used as the game specification for
    the 16-card deck and card effects. No third-party implementation was used.
\end{itemize}

\section{Conclusion}

The project delivers a coherent Love Letter AI platform with a hidden-state
environment, multiple agent families, fair and oracle search baselines, and
reproducible evaluation scripts. The strongest course connection is the
combination of partial observability, belief modeling, adversarial search, and
reinforcement learning. To reach a high final grade, the main repair is not to
add another complicated agent, but to strengthen the evidence: run a larger
symmetric tournament, report confidence intervals, include ablations, and make
the hidden-information fairness distinction explicit in both the report and
presentation.

\begin{thebibliography}{9}
\bibitem{loveletter}
Seiji Kanai. \emph{Love Letter}. Alderac Entertainment Group, base game rules.

\bibitem{pytest}
pytest documentation. \url{https://docs.pytest.org/}

\bibitem{tectonic}
Tectonic Typesetting System. \url{https://tectonic-typesetting.github.io/}
\end{thebibliography}

\end{document}
